\documentclass[english]{letter}
\usepackage[T1]{fontenc}
\usepackage[latin1]{inputenc}
\usepackage{babel}
\usepackage[top=.5in,bottom=1in,left=1in,right=1in,noheadfoot]{geometry}
\usepackage{graphicx}

\address{\sffamily
  \begin{tabular}{rr}
    Prof.~Dr.~Ant�nio Augusto Fr�hlich & \raisebox{-0.9cm}[0pt][0pt]{\resizebox*{!}{1.2cm}{\includegraphics{/home/guto/etc/logos/lisha}}}\\
    Federal University of Santa Catarina &\\
    Software/Hardware Integration Laboratory &\\
    \multicolumn{2}{r}{PO Box 476, 88040-900 Florian�polis, SC, Brazil}\\
    \multicolumn{2}{r}{Phone: +55 48 3721-9516}\\
    \multicolumn{2}{r}{E-Mail: guto@lisha.ufsc.br}\\
    \multicolumn{2}{r}{WWW: http://lisha.ufsc.br/$\sim$guto}\\
    \\
    \multicolumn{2}{r}{\today}
  \end{tabular}
}

\date{}

\signature{
  \begin{center}
   \raisebox{-2em}[0pt][0pt]{\includegraphics[scale=.75]{/home/guto/personal/signature}}\\
   \hrulefill\\
    Ant�nio Augusto Fr�hlich
  \end{center}
}


\begin{document}

\begin{letter}{International Journal of Distributed Sensor Networks\\
	Editor\\
    Prof. Cem Ersoy
}
\opening{Dear Prof. Ersoy,}
\setlength{\parskip}{1.5ex}

A revised version of the manuscript "A Comprehensive Approach to Power
Management in Embedded Systems", initially submitted to the International
Journal of Distributed Sensor Networks on September 2010, has just been uploaded
through Hindawi manuscript system. Bellow you can find replies to each
individual reviewer and also to your considerations.

\textbf{\large Reviewer 1}

\textbf{Remarks}

This paper does not suit IJDSN. The methods and tools can be applied to any
embedded system. Therefore I strongly encourage the authors to submit the paper
to a journal on embedded systems like Journal of Embedded Systems.

\textbf{Reply}

Certainly wireless sensors are not the only kind of embedded system operated on
batteries that could benefit from the power management strategy proposed in the
paper. However, from my point of view the fact that an article published on the
IJDSN addresses a problem that reaches out not only for wireless sensor networks
but for a broader range of battery-operate devices has the positive potential to
attract new readers. Indeed, the sensor/actuator system used as a didactic
example throughout the paper already establishes a bound to WSN. The strong
emphasis on energy consumption by I/O devices, particularly by radio
transceivers and analog/digital converters, would also probably be taken by
readers as WSN-focused matter. Nonetheless, two new WSN case studies in the
realm of EPOSMote have been added to Section V: the C-MAC protocol (Section V.D)
and the AD-ZRP routing protocol (Section V.E).

\textbf{\large Reviewer 2}

\textbf{Remarks}

The paper presents an approach for the design of embedded systems with both,
real-time and energy constraints. The approach comprises a design-time component
as well as a runtime autonomic management component. The focus of the integrated
approach is on the usability of power management features of peripheral
components under the precondition of real-time requirements. A special tool
modifies the function calls to accesses to peripheral devices at compile-time to
control the power managment features of the corresponding device, especially to
enable the device if necessary. Switching to power save mode must be performed
explicitly by the developer (with his knowledge of future usage) and/or in
combination with a modified software scheduler and an autonomic power manager.
The modified scheduler is aware of real-time constraints as well as the battery
utilisation and controls the power management accordingly. The approach is
demonstarted and evaluated by a single example that is explained clearly. The
results show a significant saving of energy of the standard configuration and
also show some alternative configurations with additional threads resulting in a
higher enrgy consumption and/or an impossible schedule because of the real-time
constraints. In Section V.B. Example Application, first paragraph, the authors
state that the EPOS tools deliver a "run-time library that realizes the required
interfaces and hardware description that could be matched by virtually any
hardware platform". What does this mean? Does the single run-time library
support any kind of hardware device? Or does it mean that it is possible to
generate a specialized run-time libraries for each device? Please clarify this.
A further related work that should be cited is the work from Kluge et.al:
"Optimisation of Energy Consumption of Soft Real-Time Applications by Workload
Prediction", Florian Kluge, Sascha Uhrig, J�rg Mische, Benjamin Satzger, Theo
Ungerer,13th IEEE International Symposium on Object/Component/Service-Oriented
Real-Time Distributed Computing Workshops (First IEEE Workshop on
Self-Organizing Real-Time Systems, SORT 2010), Carmona, Spain. They presented a
technique that automatically detects the schedule within an embedded system and
controls the DVS of the processor automatically. The proposed technique focuses
on soft real-time systems only. Due to some typos, proof-reading should be a
good idea. Especially, Figure, Section, Table should be written with capital
letters.

\textbf{Reply}

Thank you for your review. Your question about how EPOS tools deliver a tailored
run-time system is very pertinent. We are currently half the way in the
transition to a new tool set that relies on evolutionary algorithms to optimize
the configuration and shall be able to point out the most suitable platforms for
a given application. This led me to the confusing sentence you identified. It has
been promptly adjusted to reflect the current tool-chain, in which a two-phase
process is performed. The tools are first speculatively deployed in order to
produce an initial specification. This specification can either be taken or
manually modified. A second invocation produces the useful run-time system. The
process has been described in earlier publications, including [31]. The suggested
related work was cited on the paragraph discussing DVS and the capitalization of
text element names was done.

\textbf{\large Reviewer 4}

\textbf{Remarks}

This paper proposed a power management scheme that was designed and implemented
from user level API to the infrastructure for embedded systems. The author has
been studied power management issues through the EPOS (Embedded Parallel
Operating system) project and published several papers related to this paper.
The efficient power management is one of important issue in embedded system and
it is well structured. However I will give some comments and some problems. -
One of main contributions is to adapt the heuristic knowledge. How to present
the knowhow of designer? Does it depend on only Petri Nets? - Semantic power
modes of the proposed power management API are ambiguous to reader and user. It
is possible that user selects wrong API. How to solve this problem? - This paper
suggested that the proposed power management scheme considers real-time
features. However, the author just mentioned some strategies. What components
are designed and implemented in proposed infrastructure?

\textbf{Reply}

Thank you for your review. Capturing designers' knowledge is certainly a big
open question for Software Engineering and many other disciplines. In the paper,
I do not propose means, but strengthen that such knowledge should not be ignored
as happens with operating systems designed for end users. In the first paragraph
of Section II, traditional UML diagrams are pointed as possible means. On the
diagrams throughout the text, power management knowledge is expressed through
explicit invocations of Power() methods and also by tagging operations with
remark notes. Indeed, the Petri Nets used to model the Operation Mode Transition
Networks described in Section III.B formally capture the operating system
designer knowledge about how mode transitions affect distinct components, but
this is only one aspect of the system. Deploying it broadly would be
unpractical. The paragraph that introduces the semantic power modes in Section
II.B has been rewritten to include examples that make their role easier to
understand. The question about the implementation of power management
infrastructure components that preserve the real-time characteristics of the
run-time system in which they are to be deployed is very relevant, but
unfortunately cannot be answered in this paper. Our long-term work with static
metaprogramming and aspect-oriented programming has shown that, even if all
algorithms within a component are deterministic and can have their WCET
estimated, the adaptation of components performed by the associated compilers
and tools while generating the final library or kernel, as long as the
interference caused by one component on others, can compromise the system as a
whole. For now, this verification is empirical in EPOS.

\clearpage

\textbf{\large Editor}

\textbf{Remarks}

Although two of the referees found the paper more useful, one of the
reviewers commented that "The methods and tools can be applied to any embedded
system." Therefore, the paper is more suitable for a journal on
embedded systems like Journal of Embedded Systems. I tend to agree with that
reviewer. For that reason, I will request that you include more examples
specific to sensor networks showing how your proposed ideas will be used and
will be effective in practice. Some of the indicators showing that the paper is
written generically for embedded systems can be found on Page 4 in the parts
mentioning the file system. It sound like the envisioned platform is much more
complex than a simple sensor node. Moreover, only a few references are related
to WSN and majority of them are related to generic embedded systems. If that is
the case, I will also recommend that you submit your work to another journal
related to embedded systems.

One additional problem not mentioned by the reviewers is the fact that there is
only one simple experiment reported in Section V.B, and many conclusions are
made based on this demonstrative example. I believe that more experiments are
needed in order to be conclusive.

\textbf{Reply}

Thank you for your constructive suggestions. Since I received your e-mail, I
have been intensively working on two new experiments aiming at further
supporting my conclusions. The C-MAC protocol (Section V.D) has been conditioned
in an experiment that corroborates the proposed autonomous power manager, while
the AD-ZRP routing protocol (Section V.E) was extended to be energy-aware and
further confirmed all the proposed mechanisms.

Sincerely,

\closing{~}

\end{letter}

\end{document}





















































